\subsection*{Reading the role named in the question}
On a new song credit, the visual reader named the composer correctly but called him the producer. The displayed question explicitly asked for the composer. A classifier given the pictured question raised the number of correct relations on the 48 human questions from 42 to 46. With the existing sentence patterns, complete answers rose from 33 to 36. In a matched fit with permuted relation labels they fell to 31. All three gains were on the 29 simplified Chinese questions. The count on the 19 traditional Chinese questions stayed at 16. On independently written programme notes, complete answers rose by one to 17 of 32. On 1,200 held-out DuIE passages with generated questions, complete answers rose by 89 to 930 (Fig.~\ref{fig:pictured-relation}).

Running the autoregressive writer in both requested word orders gave two complete, stopped answers on 34 of the 48 human questions, 18 of the 32 programme questions and 892 of the 1,200 DuIE passages. Six of 2,400 individual DuIE replies reached the length limit. Both replies to the original song question now correctly identified the composer. On a separate song, three of four questions passed in both orders. After exchanging the composer and lyricist, the reader continued to credit the former composer in one answer. All four replies were correct for only two of eight exchanged programme notes with the induced writer. Recognizing a role in the question can repair a wrong attribution while leaving the choice of person unresolved.
